% Generated draft: requires mathematical and source review.
\documentclass[11pt]{article}
\usepackage[margin=25mm]{geometry}
\usepackage{fontspec}
\setmainfont{Latin Modern Roman}
\usepackage{amsmath,amssymb,mathtools,unicode-math}
\setmathfont{Latin Modern Math}
\usepackage{enumitem,needspace,xurl,hyperref}
\setlength{\parindent}{0pt}
\setlength{\parskip}{5pt}
\setlength{\emergencystretch}{3em}
\newcommand{\N}{\mathbb{N}}
\newcommand{\Z}{\mathbb{Z}}
\newcommand{\Q}{\mathbb{Q}}
\newcommand{\R}{\mathbb{R}}
\newcommand{\C}{\mathbb{C}}
\newcommand{\F}{\mathbb{F}}
\newcommand{\A}{\mathbb{A}}
\newcommand{\PP}{\mathbb{P}}
\newcommand{\E}{\mathbb{E}}
\newcommand{\eps}{\varepsilon}
\newcommand{\dd}{\,\mathrm{d}}
\begin{document}
\section*{UnsolvedMath ID: 20002463}
\textbf{Source problem code:} catalogue entry
\subsection{Definitions and mathematical statement}
A.10, ``Jin, Shi,'' is a participant-interest paragraph, not an open problem. It describes numerical methods for multiscale physics, transitions from quantum to classical mechanics and from kinetic theory to hydrodynamics, and numerical Liouville equations with singular Hamiltonians arising from discontinuous potentials or discontinuous local wave speeds. The trailing $8$ is a page marker, and ``Hamiltoni-ans'' is a line-break artifact. No equation, interface condition, numerical scheme, or requested theorem is supplied.

The aid's separate contextual lemma concerns one-dimensional scattering at an interface. On each side take $H_\pm(p)=T_\pm(p)+V_\pm$, with even kinetic energies strictly increasing in $|p|$. On transmitting positive-momentum branches assume sufficient differentiability and nonzero velocities $v_\pm(p)=T_\pm'(p)>0$. Let the energy map $P$ satisfy
\[
H_+(P(p))=H_-(p),\qquad P(p)>0.
\]
Then its Jacobian is
\[
P'(p)=\frac{v_-(p)}{v_+(P(p))}.
\]
For incident phase density $f_-$ the natural positive flux measure is $\dd\mu_-(p)=|v_-(p)|f_-(p)\,\dd p$. A conservative transmitting interface law is
\[
\mu_+=P_\#\mu_-,
\]
where the pushforward means $\mu_+(E)=\mu_-(P^{-1}(E))$. On the regular branch this is exactly the cancellation of the velocity ratio by the momentum Jacobian; for full deterministic transmission the corresponding phase densities satisfy $f_+(P(p))=f_-(p)$. It is not a rule obtained by pushing forward $f(p)\,\dd p$ alone.

More generally, energy-preserving reflected and transmitted branches may split the incident flux, with nonnegative branch weights summing to one. For every measurable weight $W$ for which the integrals exist, the reported identity is
\[
\sum_{\text{outgoing branches}}\int W(H_{\mathrm{out}}(p))\,\dd\mu_{\mathrm{out}}(p)
=\int W(H_{\mathrm{in}}(p))\,\dd\mu_{\mathrm{in}}(p).
\]
For even kinetic energy, reflection on the same side uses $p\mapsto-p$ and preserves energy. The identity covers Hamiltonian-weighted flux moments, not arbitrary momentum moments across a potential step.

The threshold qualification is essential. Suppose an uphill threshold incident momentum $p_*$ has positive incident speed, energy $H_-(p_*)=H_+(0)$, and nonzero atomic flux $\mu_-(\{p_*\})=m>0$. If $v_+(0)=0$, transmission would demand outgoing flux $m\delta_0$. But for any finite outgoing momentum occupation measure $\nu_+$,
\[
\mu_+=|v_+|\nu_+\quad\Longrightarrow\quad\mu_+(\{0\})=0.
\]
Thus that atom cannot be represented by a finite zero-speed outgoing occupation measure. A conservative model must instead reflect it or explicitly track interface storage and its flux balance. This obstruction concerns atomic threshold flux; it is not automatically present for ordinary integrable densities assigning zero mass to the threshold. The source paragraph does not specify which interface convention is intended.
\subsection{Short English statement}
Shi Jin's contribution describes research interests rather than posing a problem. The accompanying aid studies an energy-preserving interface law for Liouville transport: the conserved object is velocity-weighted flux measure. An incident atom at a transmission threshold needs reflection or interface storage because zero outgoing speed cannot carry that flux through a finite occupation measure.
\subsection{Sources}
\begin{itemize}
\item Local UnsolvedMath registry, permanent ID $20002463$, \texttt{catalogue entry}; \texttt{aim-pdes-notes.json}, \texttt{notes[74]}, section tag.
\item American Institute of Mathematics, \emph{Stiff sources and numerical methods for conservation laws}, A.10, ``Jin, Shi,'' \url{https://aimath.org/WWN/balancelaws/balancelaws.pdf}. No proposer, specific mathematical problem, or authoritative separate citation for the scattering lemma is supplied.
\item Research attempt 1: context only, status ``not a problem'' at run, rigor label ``proved,'' novelty confidence low. Registry labels open and open triage, statement status exact. The flux lemma is a synthesis prompted by the paragraph, not an original AIM problem statement.
\end{itemize}
\end{document}
